%%%%%%%%%%%%%%%%%%%%%%%%%%%%%%%%%%%%%%%%%%%%%%%%%%%%%%%%%%%%%%%%%%%%%%
% Lớp 11 | Quan hệ vuông góc trong không gian | Hai mặt phẳng vuông góc
% Dạng 1H8?4-1: Câu hỏi lý thuyết
% Số câu: 9  (sinh tự động bởi tach_cau.py; từ nay thầy sửa trực tiếp file này)
%%%%%%%%%%%%%%%%%%%%%%%%%%%%%%%%%%%%%%%%%%%%%%%%%%%%%%%%%%%%%%%%%%%%%%

%%%%%%%%%%%%%%%% TRẮC NGHIỆM 4 PHƯƠNG ÁN %%%%%%%%%%%%%%%%

% nguồn: Ok_TracNghiemLop11_Lan1
\begin{ex}%[1H8N4-1]%[ma:913ef3d7]
	Gọi $\alpha$ là góc giữa hai mặt phẳng $(P)$, $(Q)$. Mệnh đề nào đúng khi nói về  số đo của góc $\alpha$.
	 \choice
		 {$0^\circ$}
		{\True $0^\circ\le\alpha\le  90^\circ$}
		 {$90^\circ\le\alpha\le 180^\circ$}
		{$0^\circ\le\alpha\le 180^\circ$} 
	 \loigiai{
	$\alpha$ là số đo giữa hai mặt phẳng nên  $0^\circ\le\alpha\le 90^\circ$.}
\end{ex}

% nguồn: Ok_TracNghiemLop11_Lan1
\begin{ex}%[1H8N4-1]%[ma:4e120353]
	Trong không gian,  khẳng định nào sau đây \textbf{sai}?
	 \choice
		{Nếu ba mặt phẳng cắt nhau theo ba giao tuyến phân biệt thì ba giao tuyến ấy hoặc đồng  quy hoặc đôi một song song với nhau}
		{\True Hai đường thẳng phân biệt cùng vuông góc với một đường thẳng thì song song  với nhau}
		{Hai mặt phẳng phân biệt cùng vuông góc với một đường thẳng thì song song với nhau} 
		{Cho hai đường thẳng  chéo nhau. Có duy nhất một mặt phẳng chứa đường thẳng này và song song với đường thẳng kia}
	 \loigiai{
		\lq\lq  Hai đường thẳng phân biệt cùng vuông góc với một đường thẳng thì song song với nhau\rq\rq \space sai vì \lq\lq Hai đường thẳng phân biệt cùng vuông góc với một mặt phẳng thì song song với nhau\rq\rq.
	} 
\end{ex}

% nguồn: Ok_TracNghiemLop11_Lan1
\begin{ex}%[1H8N4-1]%[ma:1c1e35cf]
	Mệnh đề nào sau đây  là \textbf{đúng}?
	\choice 
		{Hai mặt phẳng vuông góc với nhau  thì mọi đường thẳng nằm trong mặt phẳng này sẽ vuông góc với mặt phẳng kia}
		{Hai mặt phẳng phân biệt cùng vuông góc  với một mặt phẳng thì vuông góc với nhau}
		{Hai mặt phẳng phân biệt cùng vuông góc  với một mặt phẳng thì song song với nhau}
		{\True Hai mặt phẳng vuông góc với nhau thì mọi đường thẳng nằm trong mặt phẳng này và vuông góc với giao tuyến của hai mặt phẳng sẽ  vuông góc với mặt phẳng kia}
	\loigiai{ 
		Mệnh đề \lq\lq Hai mặt phẳng vuông góc với nhau thì mọi đường thẳng nằm trong  mặt phẳng này sẽ vuông góc với mặt phẳng kia\rq\rq\,sai vì có thể xảy ra trường hợp hai mặt phẳng vuông góc với nhau nhưng đường thẳng thuộc mặt phẳng này song song với mặt phẳng kia.\\
		Mệnh đề \lq\lq Hai mặt phẳng phân biệt cùng vuông góc với một mặt phẳng thì vuông góc với nhau\rq\rq\,sai vì xảy ra trường hợp hai mặt phẳng song song.\\ 
		Mệnh đề \lq\lq Hai mặt phẳng phân  biệt cùng vuông góc với một mặt phẳng thì song song với nhau\rq\rq\,sai vì xảy ra trường hợp hai mặt phẳng vuông góc.
	 }
\end{ex}

% nguồn: Ok_TracNghiemLop11_Lan1
\begin{ex}%[1H8N4-1]%[ma:525ba7f1]
	Điền vào chỗ trống sau để có khẳng định đúng.\\ 
	\lq\lq Nếu hai mặt phẳng $\ldots$ $\ldots$ với nhau thì bất kì một  đường thẳng nào nằm trong mặt phẳng này và vuông góc với giao tuyến cũng vuông góc với mặt phẳng kia\rq\rq
	 \choice
		{song  song}
		{cắt nhau} 
		{phân biệt} 
		{\True vuông góc} 
	 \loigiai{
		Mệnh đề đúng \lq\lq Nếu hai  mặt phẳng \textit{\textbf{\underline{vuông góc}}} với nhau thì bất kì một đường thẳng nào nằm trong mặt phẳng này và vuông góc với giao tuyến cũng vuông góc với mặt phẳng kia\rq\rq
	 }
\end{ex}

% nguồn: Ok_TracNghiemLop11_Lan1
\begin{ex}%[1H8N4-1]%[ma:af9ab92e]
	Mệnh đề nào  sau đây đúng?
	 \choice
		{Đường thẳng $d$ là đường vuông góc chung của hai đường thẳng nhau $a$, $b$ khi và chỉ khi $d$  vuông góc với cả $a$ và $b$}
		{\True Cho hai mặt phẳng vuông góc với nhau, nếu một đường thẳng nằm trong mặt phẳng này và vuông góc với giao  tuyến của hai mặt phẳng thì vuông góc với mặt phẳng kia}
		{Nếu đường thẳng $a$ và mặt phẳng $(P)$ cùng vuông góc với đường thẳng $d$ thì đường thẳng  $a$ song song với mặt phẳng $(P)$}
		{Nếu hai  đường thẳng phân biệt cùng vuông góc với một đường thẳng thứ ba thì chúng song song với nhau}
	\loigiai{ 
	Hai  mặt phẳng vuông góc với nhau, nếu một đường thẳng nằm trong mặt phẳng này và vuông góc với giao tuyến của hai mặt phẳng thì vuông góc với mặt phẳng kia.}
\end{ex}

% nguồn: Ok_TracNghiemLop11_Lan1
\begin{ex}%[1H8N4-1]%[ma:6e595ed0]
	Trong các mệnh đề sau,  mệnh đề nào là \textbf{đúng}?
	 \choice
		{Nếu hai mặt phẳng vuông góc với nhau thì mọi đường thẳng nằm trong  mặt phẳng này sẽ
		 vuông góc với mặt phẳng kia}
		{Hai mặt phẳng phân biệt cùng vuông  góc với một mặt phẳng thì vuông góc với nhau}
		{\True Hai mặt phẳng vuông góc với nhau thì mọi đường thẳng nằm trong mặt phẳng  này và vuông
		góc với giao tuyến của hai  mặt phẳng sẽ vuông góc với mặt phẳng kia}
		 {Hai mặt phẳng phân biệt cùng vuông góc với một mặt phẳng thì song song với nhau}
	 \loigiai{
	 }
\end{ex}

% nguồn: Ok_TracNghiemLop11_Lan1
\begin{ex}%[1H8N4-1]%[ma:f95d4645]
	Khẳng định nào  sau đây \textbf{đúng}?
	 \choice
		{\True Hình lăng trụ  đứng có đáy là một đa giác đều là hình lăng trụ đều}
		{Hình lăng trụ đứng là  hình lăng trụ đều}
		{Hình lăng trụ có đáy là một đa giác đều là hình  lăng trụ đều}
		{Hình lăng  trụ tứ giác đều là hình lập phương}
	 \loigiai{
		 Hình lăng trụ đứng có đáy là một đa giác đều là hình lăng trụ đều.
	 }
\end{ex}


%%%%%%%%%%%%%%%% ĐÚNG/SAI %%%%%%%%%%%%%%%%

% nguồn: Ok_dungsai_lop11_Lan1
\begin{ex}%[1H8N4-1]%[ma:cba44de7]
 Cho hai mặt phẳng $(P)$ và $(Q)$  cắt nhau theo giao tuyến $a$.
 \choiceTF 
   {\True $(P)$ vuông góc với $(Q)$ nếu góc giữa chúng bằng $90^{\circ}$}
  {Nếu $\alpha $ là  góc giữa $(P)$ và $(Q)$ thì $0^{\circ}<\alpha <180^{\circ}$}
  {\True Nếu $(P)$ chứa đường thẳng vuông góc với $(Q)$ thì $(P)$ cũng vuông  góc với $(Q)$}
  {\True Nếu $(P)$ và $(Q)$ vuông góc với nhau thì bất kì đường thẳng nào nằm trong  $(P)$ mà vuông góc với $a$ cũng vuông góc với $(Q)$}
  \loigiai{
  \begin{itemchoice}
   \itemch Đúng.  Theo định nghĩa hai mặt phẳng vuông góc.
   \itemch Sai. Nếu $\alpha  $ là góc giữa $(P)$ và $(Q)$ thì $0^{\circ}<\alpha\le{90^{\circ}}$.
   \itemch Đúng. Theo định  lí hai mặt phẳng vuông góc.
   \itemch Đúng. Theo hệ  quả của định lý hai mặt phẳng vuông góc.
  \end{itemchoice}
 } 
\end{ex}

% nguồn: Ok_dungsai_lop11_Lan1
\begin{ex}%[1H8B4-1]%[ma:0cf968ee]
 Xét các khẳng định  nào sau
  \choiceTF
  {\True $\heva{&(\alpha) \cap (\beta) = d\\& (\alpha) \perp (P) \\& (\beta) \perp (P)}\Rightarrow d \perp (P)$} 
  {\True $\heva{&(\alpha)\parallel (\beta) \\&  ( P) \perp ( \alpha)}
  \Rightarrow ( P)  \perp ( \beta)$}
  {$\heva{&(\alpha) \perp (\beta)  \\& a \subset (\alpha) \\& b \subset ( \beta)} \Rightarrow a \perp b$}
  {$\heva{&(\alpha) \perp (\beta) \\&  a \subset (\alpha)} \Rightarrow a \perp (\beta)$}
 \loigiai{ 
  \begin{center}
   \begin{tikzpicture}[scale=0.7,font=\footnotesize,line join=round,line cap=round,>=stealth]
     \path 
       (0,0) coordinate (B)
      (1,1.4) coordinate  (A)
      (4,0) coordinate  (C)
       ($(C)-(B)+(A)$) coordinate (D)
      ($(A)+(90:3.5)$)  coordinate (A')
       ($(B)-(A)+(A')$) coordinate (B')
       ($(C)-(A)+(A')$) coordinate (C')
      ($(D)-(A)+(A')$)  coordinate (D');
     \draw  (B')--(B)--(C)--(D)--(D')--(A')--(B')--(C')--(D')--cycle (C')--(C)
      (B')--(C); 
     \draw[dashed] (A')--(A)--(B) 
      (A)--(D); 
     \foreach \p/\q in {A/160,B/-135,C/-45,D/0,A'/90,B'/180,C'/-20,D'/90} 
     \fill[black] (\p) node[shift={(\q:3mm)}]{$\p$} circle (1.0pt); 
    \end{tikzpicture}
  \end{center}
  \begin{itemchoice}
   \itemch Đúng, vì hai  mặt phẳng cùng vuông góc với mặt phẳng thứ 3 thì giao tuyến của chúng cũng vuông góc.
   \itemch Đúng, vì một mặt phẳng vuông góc với một trong hai mặt phẳng song song thì sẽ vuông góc với mặt  phẳng còn lại.
   \itemch Sai vì, xét hỉnh lập phương $ABCD.A'B'C'D'$  có
   \[\heva{& (A'B'C'D') \perp (BCC'B') \\& B'D'  \subset (A'B'C'D'), B'C \subset (BCC'B') \\ & BC', B'D' \text{ không vuông góc}}\]
   \itemch Sai vì, xét hỉnh lập  phương $ABCD.A'B'C'D'$ có
   \[\heva{& (A'B'C'D') \perp (BCC'B') \\& B'D' \subset (A'B'C'D') \\ &  B'D' \text{ không vuông góc với } (BCC'B')}\]
  \end{itemchoice}
  }
\end{ex}
